\scanpage{8}

\begin{center}
U N I V E R S I T E\\
DES SCIENCES ET TECHNIQUES\\
\uline{DU LANGUEDOC}

Service de la Scolarité\\
MA/GJ\\
Diffusion : Liste A\\
N° : 4679
\end{center}

Montpellier, le 20 octobre 1983

\begin{center}
\uline{AVIS DE SOUTENANCE D'UNE THESE DE DOCTEUR INGENIEUR}

\uline{AUTOMATIQUE}
\end{center}

\uline{Nom et prénom du candidat} : SITOHANG Benhard

\uline{Titre de la thèse} : LINAIRE : SYSTEME DE GESTION DE BASES DE DONNEES RELATIONNELLES SUR MICRO-ORDINATEUR.

\uline{Date, heure et lieu de soutenance prévus:} 28 OCTOBRE 1983 à 14h30 - CRIG USTL

\uline{Nom du Directeur de thèse} : Monsieur le Professeur B. FILLIATRE

\uline{Désignation du laboratoire où la thèse est préparée} : CRIG USTL

\uline{Résumé de la thèse} :

La conception et la réalisation d'un Système de Gestion de Bases de Données Relationnelles sur micro-ordinateur, que nous avons baptisé ``LINAIRE'', est le sujet de cette thèse.

Nous mettons l'accent sur trois points principaux :

1. BOUCLE MULTI-NIVEAUX ; c'est une technique du traitement de données, permettant d'éviter l'utilisation des relations intermédiaires comme résultats des opérateurs relationnels, ce qui serait particulièrement pénalisant pour un micro-ordinateur possédant une mémoire centrale limitée.

2. STRUCTURE INTERNE DE LA BASE ; elle est multiple, afin de préserver au mieux les performances tout en offrant en particulier la possibilité d'accélérer le traitement de l'opérateur relationnel ``équi-joint''.

3. L'adaptation sur un micro-ordinateur du langage de manipulation de QUERY-BY-EXAMPLE (langage de haut niveau, interactif, et tabulaire).

\uline{Durée approximative des recherches ayant abouti à la thèse} : 2 ans

\uline{Composition du jury proposé} :

\begin{tabular}{ll}
MM. & FILLIATRE B.\\
& CHEIN M.\\
& PICHAT E.\\
& PIROTTE A.\\
& PRUNET F.
\end{tabular}

\begin{center}
P. le Secrétaire Général,\\
l'Attachée Principale d'Administration Scolaire\\
et Universitaire,

\unsure{A. Gentil}

A. GENTIL
\end{center}

\margin{7}
